\section*{Аннотация}
Данная работа посвящена исследованию улучшения точности прогнозирования хаотических временных 
рядов путем расширения обучающей выборки данными из схожих динамических систем. В ходе исследования
 на данных системы Лоренца было установлено, что не всякая формально близкая система является
  полезным донором данных, и их неверный выбор может ухудшить качество прогноза. Для решения 
  этой проблемы был разработан критерий для априорного отбора совместимых
   рядов, основанный на сравнении их статистических инвариантов. Эксперименты подтвердили, что
    использование этой метрики позволяет целенаправленно улучшать прогноз, добиваясь значительного
     сокращения непрогнозируемых точек при сохранении высокой точности.

\subsection*{Ключевые слова}
Прогнозирование хаотических временных рядов, система Лоренца,
расширение обучающей выборки, метод мотивов, непрогнозируемые точки